% Part: propositional-logic
% Chapter: propositional-logic
% Section: preliminaries

\documentclass[../../../include/open-logic-section]{subfiles}

\begin{document}

\olfileid{pl}{syn}{pre}

\olsection{પ્રારંભિક બાબતો}

\begin{thm}[\emph{!!{formula}s પર અનુમાનનો સિદ્ધાંત}]
  \ollabel{thm:induction}
  કોઈ ગુણધર્મ~$P$ બધાં આણ્વિક !!{formula}s માટે હોય અને નીચેની
  શરતો સંતોષતો હોય:
  \begin{enumerate}
    \tagitem{prvNot}{$!A$ માટે ગુણધર્મ હોય ત્યારે $\lnot !A$ માટે
      પણ હોય;}{}
    \tagitem{prvAnd}{$!A$ અને~$!B$ માટે ગુણધર્મ હોય ત્યારે
      $(!A \land !B)$ માટે પણ હોય;}{}
    \tagitem{prvOr}{$!A$ અને~$!B$ માટે ગુણધર્મ હોય ત્યારે
      $(!A \lor !B)$ માટે પણ હોય;}{}
    \tagitem{prvIf}{$!A$ અને~$!B$ માટે ગુણધર્મ હોય ત્યારે
      $(!A \lif !B)$ માટે પણ હોય;}{}
    \tagitem{prvIff}{$!A$ અને~$!B$ માટે ગુણધર્મ હોય ત્યારે
      $(!A \liff !B)$ માટે પણ હોય;}{}
  \end{enumerate}
  તો બધાં !!{formula}s માટે $P$ હોય છે.
\end{thm}

\begin{proof}
  ગુણધર્મ~$P$ ધરાવતાં બધાં !!{formula}sનો સંગ્રહ $S$ લો. સ્પષ્ટપણે
  $S \subseteq \Frm[L_0]$. $S$, \olref[fml]{defn:formulas}ની બધી
  શરતો સંતોષે છે: તેમાં બધાં આણ્વિક !!{formula}s છે અને તે
  !!{operator}s હેઠળ સંવૃત છે. $\Frm[L_0]$ આવો નાનામાં નાનો વર્ગ છે,
  તેથી $\Frm[L_0] \subseteq S$. આમ $\Frm[L_0]=S$, અને દરેક !!{formula}ને
  ગુણધર્મ~$P$ છે.
\end{proof}

\begin{prop}\ollabel{prop:balanced}
   $\Frm[L_0]$નું દરેક !!{formula} \emph{સંતુલિત} છે, એટલે તેમાં
   ડાબાં અને જમણાં કૌંસની સંખ્યા સમાન છે.
\end{prop}

\begin{prob}
\olref[pl][syn][pre]{prop:balanced} સાબિત કરો.
\end{prob}

\begin{prop} \ollabel{prop:noinit}
!!a{formula}નો કોઈ ઉચિત આરંભખંડ !!a{formula} નથી.
\end{prop}

\begin{prob}
  \olref[pl][syn][pre]{prop:noinit} સાબિત કરો.
\end{prob}

\begin{prop}[એકમાત્ર વાચનીયતા]
$\Frm[L_0]$નું દરેક !!{formula}~$!A$ નીચેના પૈકી કોઈ એક સ્વરૂપમાં
બરાબર એક જ રીતે વિભાજિત થઈને વાંચી શકાય છે:
\begin{enumerate}
\tagitem{prvFalse}{$\lfalse$.}{}

\tagitem{prvTrue}{$\ltrue$.}{}

\item કોઈ $\Obj p_n \in \PVar$ માટે $\Obj p_n$.

\tagitem{prvNot}{કોઈ !!{formula}~$!B$ માટે $\lnot !B$.}{}

\tagitem{prvAnd}{કોઈ !!{formula}s $!B$ અને~$!C$ માટે
$(!B \land !C)$.}{}

\tagitem{prvOr}{કોઈ !!{formula}s $!B$ અને~$!C$ માટે
$(!B \lor !C)$.}{}

\tagitem{prvIf}{કોઈ !!{formula}s $!B$ અને~$!C$ માટે
$(!B \lif !C)$.}{}

\tagitem{prvIff}{કોઈ !!{formula}s $!B$ અને~$!C$ માટે
$(!B \liff !C)$.}{}
\end{enumerate}
વળી, આ વિભાજિત વાંચન \emph{એકમાત્ર} છે.
\end{prop}

\begin{proof}
$!A$ પર અનુમાનથી. દાખલા તરીકે, ધારો કે $!A$નાં $(!B \lif !C)$ અને
$(!B' \lif !C')$ તરીકે બે જુદાં વાંચન છે. ત્યારે $!B$ અને $!B'$ સરખાં
હોવાં જોઈએ (નહિતર એક બીજાનો ઉચિત આરંભખંડ હોત); તેથી $!A$નાં બંને
વાંચન જુદાં હોય તો એનું કારણ માત્ર એ હોઈ શકે કે $!C$ અને $!C'$ એક જ
સંકેતશ્રેણીનાં જુદાં વાંચન છે, જે અનુમાન-પૂર્વધારણાથી અશક્ય છે.
\end{proof}


\begin{defn}[એકરૂપ પ્રતિસ્થાપન]
જો $!A$ અને $!B$ !!{formula}s હોય અને $\Obj p_i$ !!{propositional
variable} હોય, તો $\Subst{!A}{!B}{\Obj p_i}$ એ $!A$માં $\Obj p_i$ના
દરેક આવર્તનને $!B$ના આવર્તનથી બદલવાથી મળતું પરિણામ દર્શાવે છે;
એવી જ રીતે, $\Obj p_1$, \dots,~$\Obj p_n$ને !!{formula}s $!B_1$,
\dots,~$!B_n$થી એકસાથે કરેલું પ્રતિસ્થાપન
$\SSubst{!A}{\subst{!B_1}{\Obj p_1},\dots,\subst{!B_n}{\Obj p_n}}$થી
દર્શાવાય છે.
\end{defn}

\begin{prob} નીચેનાં પાંચ !!{formula}sમાંના દરેક માટે નક્કી કરો કે
 તેને પ્રતિસ્થાપન \( \Subst{!A}{!B}{\Obj p_i} \) તરીકે લખી શકાય છે કે નહિ,
 જ્યાં \( !A \) અનુક્રમે (i) \( \Obj p_0 \); (ii) \( ( \lnot \Obj p_0 \land \Obj
 p_1) \); અને (iii) \( ( ( \lnot \Obj p_0 \lif \Obj p_1 ) \land \Obj
 p_2 ) \) છે. દરેક કિસ્સામાં સંબંધિત પ્રતિસ્થાપન જણાવો.
  \begin{enumerate}
    \item \( \Obj p_1 \)
    \item \( ( \lnot \Obj p_0 \land \Obj p_0 ) \)
    \item \( ( ( \Obj p_0 \lor \Obj p_1 ) \land \Obj p_2 )  \)
    \item \( \lnot ( ( \Obj p_0 \lif \Obj p_1 ) \land \Obj p_2 ) \)
    \item \( (( \lnot ( \Obj p_0 \lif \Obj p_1 ) \lif ( \Obj p_0 \lor \Obj p_1 )) \land \lnot ( \Obj p_0 \land \Obj p_1 )) \)
  \end{enumerate}
\end{prob}

\begin{prob}
અનુમાન વડે $\Subst{!A}{!B}{p}$ની ગાણિતિક રીતે ચુસ્ત વ્યાખ્યા આપો.
\end{prob}

\end{document}
